\section{What the VQC Head Does to the Backbone}
\label{sec:backbone}

\subsection{The default VQC head lowers bottleneck probe accuracy}
\label{sec:gap}

On CIFAR-10 with 4,000 images, the common probe reaches $69.8\pm1.0$ on backbones trained
with the default VQC head and $74.5\pm0.5$ on backbones trained with a linear head, a paired
difference of $-4.7\pm0.9$ points that is negative for all five seeds. The probe reads the
VQC-trained bottleneck better than the VQC head itself does ($67.8\pm1.3$ test): replacing
the VQC head by the probe raises accuracy by about 2.0 points, and a 4.7-point gap remains
between the two backbones under the common probe. The gap is larger with less data (CIFAR-10:
$-11.2\pm4.0$ and $-11.1\pm1.5$ at 1,000 and 2,000 images) and smaller on MNIST ($-7.9\pm4.2$,
$-2.2\pm0.8$ and $-1.0\pm0.5$ at 250, 500 and 1,000; 3 seeds each). A narrow classical head
lowers probe accuracy as well: MLP-8 backbones probe $3.6\pm0.6$ points below linear-head
backbones on CIFAR-10/4,000, $1.0\pm0.9$ above the VQC. The rest of this section shows that
the VQC's share of the gap depends on its CNOT structure, which no classical head has.

\subsection{Removing the CNOT rings recovers most of the gap}
\label{sec:ring}

Deleting every CNOT and keeping the rest of the head unchanged raises probe accuracy to
$73.6\pm0.5$ and test accuracy to $73.1\pm0.5$ (5 seeds). Against the default VQC this is
$+3.8\pm1.2$ probe and $+5.2\pm1.4$ test points, positive for every seed, and the ring-free
head now probes $2.8\pm0.6$ points above MLP-8. It leaves $0.8\pm0.5$ probe and $1.4\pm0.5$
test points to the linear head. MNIST shows the same direction: ring removal gains
$2.8\pm0.4$ and $1.4\pm0.4$ probe points at 500 and 1,000 images, and the ring-free head's
test accuracy (94.5, 96.4) equals the linear head's (94.3, 96.4) to within 0.2 points.

\subsection{Product order tracks the gap}
\label{sec:order}

Which property of the rings matters? \Cref{fig:hero}b places every head at the product order
of \cref{eq:identity}. Across eleven heads at four orders, every head probes below every
head of lower order. A ring after the last layer only keeps the default's parameter count and
number of rotation layers but raises the order to 5.25; it lowers probe accuracy by
$9.7\pm2.2$ points relative to the default (test $-12.5\pm2.3$; 5 seeds), and the 1- and
2-layer circuits at the same order probe at 56.8 and 61.2 (3 seeds). Across these
circuits, depth and parameter count do not order the cost; product order does.

Two interventions test the products directly. First, classical heads trained end to end on a
circuit's zero-rotation products, with no circuit, land at the circuit's order on the trend:
the products of the default circuit give $68.8\pm1.6$ probe and $67.1\pm1.3$ test, paired
differences to the default VQC of $-1.0\pm2.2$ and $-0.7\pm1.7$ with mixed signs across
seeds. Second, the
\emph{pairs} circuit keeps a CNOT after every layer but has order 1, because its pair CNOTs
cancel at zero rotations. It recovers $3.1\pm1.2$ probe and $4.3\pm1.2$ test points over the
default, positive for every seed, and sits $0.7\pm1.0$ probe points below the ring-free
circuit. Pairs after the
last layer only (order 1.5) probe at $71.2\pm0.7$, between orders 1 and 2.5, as does the
classical head on their products ($71.2\pm1.3$).

The identity is exact only at zero rotations, so we check how far the trained circuits move
from it (\cref{app:zeroangle}). A circuit frozen at its initialization reproduces its
zero-rotation model with $R^2=0.997$. Trained default and last-ring circuits stay close: a
linear map of the zero-rotation readouts recovers 87\% and 90\% of the trained readout
variance, and the input Jacobians of the logits align with cosine 0.77 and 0.80. Low-order
circuits move further ($R^2$ 0.74--0.83); for the ring-free circuit this is the
$\sin\enc{j}$ term of \cref{eq:ringfree}, which its trained rotations use. At order 5.25 the trained rotations recover part
of the cost: the last-ring circuit beats the classical head on its products by
$6.4\pm5.8$ test points. The products thus reproduce the accuracy trend and explain most of
the trained ring circuits' readout variance.

\subsection{Without CNOTs the head is a classical sinusoidal model}
\label{sec:sinusoid}

\begin{table*}[t]
\centering
\caption{Without entangling gates (and with single-qubit $Z$ readouts) the VQC head is exactly a classical per-coordinate sinusoidal model,
$\mathrm{logits} = W[\alpha_j\cos a_j + \beta_j\sin a_j]_{j=1}^{8} + b$ with $a_j = \pi\tanh z_j$
and $\alpha_j, \beta_j$ set by the trained rotations, for all rotation angles. Classical heads
$W[\cos a_j]_j + b$ (a restriction) and $W[\cos a_j, \sin a_j]_j + b$ (an unrestricted
extension), trained end to end with the same recipe and per-head validation lr selection,
trail it by 0.3--0.9 probe points and 0.6--1.3 test points.
CIFAR-10 with 4,000 training images; test accuracy of the full model and of the common linear
probe on the bottleneck (\%, mean $\pm$ SD over $n$ seeds); differences in percentage points,
paired over seeds.}
\label{tab:ringfree}
\small
\setlength{\tabcolsep}{4pt}
\begin{tabular}{lccccc}
\toprule
& & & & \multicolumn{2}{c}{$-$ ring-free VQC} \\
\cmidrule(lr){5-6}
Head & $n$ & Test & Probe & Test & Probe \\
\midrule
VQC, no CNOT rings & 5 & $73.1 \pm 0.5$ & $73.6 \pm 0.5$ & --- & --- \\
cos-only (restriction) & 5 & $72.5 \pm 0.5$ & $73.3 \pm 0.4$ & $-0.6 \pm 0.7$ & $-0.3 \pm 0.6$ \\
cos+sin (unrestricted extension)$^a$ & 3 & $72.0 \pm 0.2$ & $72.8 \pm 0.4$ & $-1.3 \pm 0.4$ & $-0.9 \pm 0.5$ \\
linear (reference) & 5 & $74.5 \pm 0.7$ & $74.5 \pm 0.5$ & $+1.4 \pm 0.5$ & $+0.8 \pm 0.5$ \\
\bottomrule
\end{tabular}

\smallskip
{\footnotesize $^a$3 seeds; differences are against the ring-free VQC on the same seeds (test 73.3, probe 73.7).}
\end{table*}


By \cref{eq:ringfree}, the ring-free head computes
$W[\alpha_j\cos\enc{j}+\beta_j\sin\enc{j}]_j+b$ for every setting of its rotations.
\Cref{tab:ringfree} compares it with classical heads of this form trained under the same
protocol. The cos-only restriction trails it by $0.3\pm0.6$ probe and $0.6\pm0.7$ test points
with mixed signs, and the unrestricted cos+sin extension by 0.9 and 1.3 (3 seeds). The match
holds at the level of the function class, not the parameters: the trained ring-free circuit
has mean absolute $R_y$ angle 0.41 rad and uses the sine component (Jacobian
cosine 0.41 against the cos-only model; \cref{tab:zeroangle}).

\subsection{Controls that do not recover the gap}
\label{sec:controls}

\begin{table}[t]
\centering
\caption{Controls on CIFAR-10 with 4,000 training images. Entries are differences in the test
accuracy of the common linear probe on the frozen 8-dimensional bottleneck, in percentage
points, paired over 3 seeds (mean $\pm$ SD). The five controls leave a substantial probe gap;
removing the CNOT rings recovers most of it.}
\label{tab:controls}
\small
\setlength{\tabcolsep}{3pt}
\begin{tabular}{>{\raggedright\arraybackslash}p{2.8cm}>{\raggedright\arraybackslash}p{4.2cm}c>{\raggedright\arraybackslash}p{4.4cm}}
\toprule
Control & Comparison & Probe $\Delta$ & Observation \\
\midrule
Reference: default VQC & VQC $-$ linear & $-4.7 \pm 1.2$ & the gap under study \\
\midrule
Backbone lr$^a$ & VQC $-$ linear, backbone lr 0.0025 & $-6.3 \pm 1.2$ & gap persists at every tested backbone lr where the VQC trains \\
 & VQC $-$ linear, backbone lr 0.005 & $-4.9 \pm 1.2$ &  \\
 & VQC $-$ linear, backbone lr 0.015 & $-8.6 \pm 0.5$ &  \\
Freeze circuit at init & frozen $-$ trained circuit & $-1.6 \pm 0.7$ & gap persists with rotations frozen at initialization \\
BatchNorm before encoding & VQC+BN $-$ VQC & $-1.2 \pm 2.2$ & pre-encoding BN does not recover the probe \\
 & VQC+BN $-$ linear & $-5.9 \pm 3.4$ &  \\
Monotone encoding$^b$ & monotone $-$ default & $+0.3 \pm 0.7$ & neither encoding recovers more than 0.7 points \\
 & symmetric control $-$ default & $+0.7 \pm 0.5$ &  \\
 & monotone $-$ symmetric control & $-0.4 \pm 0.9$ &  \\
\midrule
Contrast: remove CNOT rings & ring-free $-$ default & $+3.8 \pm 1.6$ & recovers most of the gap (\S\ref{sec:ring}) \\
 & ring-free $-$ linear & $-0.8 \pm 0.7$ &  \\
\bottomrule
\end{tabular}

\smallskip
\begin{minipage}{0.97\textwidth}\footnotesize $^a$Head lr fixed at 0.01 (VQC) and 0.03 (linear) while the backbone lr varies; the VQC collapses at the tested backbone lrs 0.03 and 0.06. Other rows: each head at its validation-selected head lr, backbone lr $0.5\times$ head lr. $^b$Monotone: $R_y(\frac{\pi}{2}\tanh z + \frac{\pi}{2})$; symmetric control: $R_y(\frac{\pi}{2}\tanh z)$, same range width; default: $R_y(\pi\tanh z)$. Not a probe difference: the final-epoch pre-clip backbone gradient norm is larger with the VQC head than with the linear head ($1.92 \pm 0.13$ vs $0.75 \pm 0.02$).\end{minipage}
\end{table}


\Cref{tab:controls} tests alternative explanations on CIFAR-10/4,000. Lowering or raising the
backbone learning rate with the head learning rate fixed leaves a gap of 4.9--8.6 points at
every rate where the VQC trains, and the VQC collapses at the two highest rates. The
backbone's gradient norm is larger, not smaller, with the VQC head. Freezing the circuit's
rotations at initialization, which keeps the head close to its zero-rotation model, does
not recover the probe ($-1.6\pm0.7$ against the trained circuit). BatchNorm before the encoding
and two alternative encodings, one of which breaks the sign symmetry of
$\cos(\pi\tanh z)$, recover at most 0.7 probe points. Lowering product order gives by far the
largest recoveries.
